\section{Discriminant comparison and uniform chambers}\label{sec:A}

We establish three comparison results for generalized Calabi--Yau data.
Their proofs reduce to bounded generic Fano fibres, but impose no bound on
the Cartier indices of the nef b-divisors.
The geometric input is a family of charts with a uniform bound on angular
variation of discrepancies; the contact of a trait with the parameter
space is allowed to be arbitrarily large.

\subsection{Conventions and the canonical bundle formula}

We use the usual generalized pair notation.
A nef part \(\mathbf M\) is a \(\mathbb Q\)-Cartier b-divisor determined and nef on some model (nef over the point unless indicated).
In particular all pullbacks involved in forming generalized discrepancies are ordinary pullbacks on sufficiently high models.
Write \(a(v)\) for log discrepancy, extended to positive multiples of \({\rm ord}_E\) by homogeneity; unless specified, numerical conditions on divisorial valuations are imposed for \emph{primitive} ones (\({\rm ord}_E\)).
Coefficient functions of b-divisors are extended by homogeneity in the same way.
Under a generically finite pullback use log pullback, so discrepancies just compose with restriction (ramification formula, characteristic zero).
Restriction of a divisorial valuation to a finitely generated subfunction
field is a multiple of a divisorial valuation or trivial.
Indeed, let \(L/K\) be the extension, of transcendence degree \(q\), and
let \(v\) be divisorial on \(L\), with nontrivial restriction \(w\).
Both value groups have rational rank one and
\[
 \trdeg_{\kappa(w)}\kappa(v)\le q,
 \qquad \trdeg_k\kappa(w)\le\trdeg_k K-1.
\]
Equality \(\trdeg_k\kappa(v)=\trdeg_k L-1\) forces equality in
both inequalities.  Thus \(w\) is divisorial.  If \(v\) has integral
values, its restricted value group is a subgroup of \(\Z\), so the
multiple of the primitive order is a positive integer.

Here and below an \emph{effective generalized CY over \(S\)} (possibly not lc) on \(X\to S\) means data as above with \(B\ge0\) on \(X\) and \(K_X+B+\mathbf M_X\sim_{\mathbb Q} f^*L\) for some \(\mathbb Q\)-Cartier \(L\) on \(S\).
CY will also indicate its discrepancy function.
Assertions about varying pairs below need only be read on the \emph{rational} parameter values (piecewise linear discrepancy functions can of course be extended on their intervals).
For an affine discrepancy family it suffices to choose common resolutions for finitely many defining members: differences of generalized discrepancies are Cartier b-evaluations.
No denominator bounds on the nef parts are assumed.
Varieties in the statements below are projective normal in characteristic zero; one can work over algebraically closed fields.
The bounds/sequence statements work also with fields allowed to vary: boundedness, complements and families used below are uniform in characteristic zero (or one can work with descents to characteristic zero subfields embeddable in \(\mathbb C\)).
We limit the dimension of the total space by a fixed integer.

We use the klt MMP with big boundary and the extraction and Mori dream
space consequences of \cite{BCHM}; preservation of Fano type under rational
contractions \cite[p.~155 and Lemma~2.8]{ProkhorovShokurov}; rational connectedness of
projective klt log Fano pairs
\cite[Corollary 1.13]{HaconMcKernan}\cite[Theorem 1]{ZhangRC};
BAB \cite[Theorem 1.1]{BirkarBAB}; and ordinary bounded lc complements
\cite[Theorems 1.7 and 1.8]{BirkarComplements}.
For complements, the seed coefficients used below are only \(0\) and
\(1\), the varieties are of Fano type over the relevant base, and the
negative log canonical divisor is nef.  The conversion to ordinary data
will be proved below; no generalized complement theorem with an
unstated nef-index hypothesis is used.
We also use the generalized klt-trivial canonical bundle formula
\cite[Theorem~1.4 and Sections~4--5]{Filipazzi}.
Recall the latter in the projective rational case: for a contraction \(g:X\to W\), data with nef \(\mathbb Q\)-Cartier b-part, \(K_X+B+\mathbf M_X\sim_{\mathbb Q}g^*L_W\), \(B\) effective over the generic point, generalized klt generically there, the discriminant is given on every model by
\begin{equation}\label{eq:A-discriminant-minimum}
   b(P)=\min_{T|_{k(W)}=e\,{\rm ord}_P,\ e>0} a(T)/e
\end{equation}
(log-discrepancy convention) and \(L_W\) minus the sum of canonical and discriminant-boundary terms (that is, \(L_W-(K_W+B_W)\), on higher models using pullback of \(L_W\)) gives a b-nef \(\mathbb Q\)-Cartier moduli part.
This is the b-nefness/b-Cartier theorem for a generalized klt-trivial fibration; the generic rank condition follows by generic klt and horizontal effectivity on \(X\) (on a resolution the rounded-up negative boundary over the generic field is effective exceptional), using \(g_*\mathcal O_X=\mathcal O_W\).
All data here are projective with nef part over the point, NQC by rationality.
We only need the generically klt case as stated here (no global klt condition).
To fit the subboundary convention in the cited theorem when some vertical
trace coefficients exceed one, subtract \(g^*G\) for a sufficiently large
effective Cartier divisor \(G\) on the base containing the images of
those components, and replace \(L_W\) by \(L_W-G\).  Generic klt is
unchanged.  The discriminant boundary changes from \(B_W\) to
\(B_W-G\), so the moduli b-divisor is unchanged.  Thus this harmless
shift permits the stated application to arbitrary vertical coefficients.
Minima are finite and attained (work over the generic DVR of \(P\)).
We apply this to \(g\) over \(S\) with \(L_W\) a pullback of \(L\).
If \(B\) is effective on \(X\), the induced trace \(B_W\) is effective because any prime on \(W\) is covered by a prime in its inverse image.
Discriminants compose by the minimum formula, including through crepant changes and birational contractions preserving the CY discrepancy data.
In fact pushing traces in small modifications and birational contractions over \(S\) (not extracting) preserves effectivity and discrepancies; crepancy follows by writing total log canonical divisor as a pullback from \(S\) plus a rational principal divisor.

\subsection{The comparison statements}

Here are three sequence versions.
Subscripts indexing a sequence will often be suppressed; ``eventually''
allows discarding finitely many indices.  Every assertion about a test
sequence of valuations is required on every subsequence of the original
sequence as well.  Bounds may depend on the fixed dimension and on the
fixed parameters explicitly appearing in a statement.  In particular,
a condition on every test sequence is uniform against a diagonal choice
of counterexamples.  A subsequent diagonal passage may make any
auxiliary divergent parameter grow more slowly.

\begin{lemma}[Comparison along long parameter intervals]\label{lem:A-comparison} Let \(X\to S\) be a sequence of Fano type contractions as above.
Suppose \(a_t\) are effective generalized CYs over \(S\), continuous finite rational piecewise affine families of discrepancies (with a finite subdivision of parameter per example, so the pieces can be treated on a common resolution), on \(0\le t\le L_i\) or \(-L_i\le t\le L_i\) with \(L_i\to\infty\).
Pieces and endpoints can be taken rational.
Values at each valuation are concave in \(t\), and \(a_0\) is klt.
Assume the corresponding one of the following (one-sided for \(0\le t\le L_i\); two-sided for \(-L_i\le t\le L_i\)):
\begin{enumerate}[label=\textup{(\roman*)}]
\item (one-sided) the discrepancy functions are nonincreasing, and whenever \(a_0(T)\to0\), \(a_t(T)/a_0(T)\to1\) for every fixed \(t\);
\item (two-sided) for some fixed \(\delta>0\), whenever \(a_0(T)<\delta\), \(a_t(T)/a_0(T)\to1\) for every fixed signed \(t\).
\end{enumerate}

As throughout these lemmas, tests on sequences of valuations include along subsequences.
Then on every fixed bounded parameter range the data are eventually klt
over the generic point of \(S\).  Write \(b_t\) for their discriminant
discrepancies on \(S\).  In case \textup{(i)}, they are nonincreasing
and satisfy \(b_t(P)/b_0(P)\to1\) for every fixed \(t\) whenever
\(b_0(P)\to0\).  In case \textup{(ii)}, some fixed \(\delta'>0\)
gives the same convergence for every fixed signed \(t\) whenever
\(b_0(P)<\delta'\).  These conclusions include every subsequence and
primitive prime orders on every birational model of \(S\).

\end{lemma}

\begin{lemma}[Phases under pure pullback]\label{lem:A-pullback-phases} Let \(a\) be a sequence of klt effective generalized CYs over \(S\) in the same Fano type setting, and let \(H\) be a sequence of \(\mathbb Q\)-Cartier \(\mathbb Q\)-divisors on the respective bases \(S\).
Suppose for some fixed positive integer \(m\), for every \(a(T)\to0\) the phase \(v_T(m f^*H)\bmod 1\) has lifts to \(\mathbb R\) of size \(O(a(T))\) bounded below by \(-o(a(T))\).
Then with discriminant \(b\), for some fixed positive multiple \(m'\), \(v_P(m'H)\) satisfies the same phase condition wherever \(b(P)\to0\).

\end{lemma}

\begin{lemma}[A uniform linear chamber]\label{lem:A-linear-chamber} In the same setting suppose \(a\) is an lc effective generalized CY over \(S\), with values in \((1/p)\mathbb Z\) for fixed \(p\), and \(a+t d\) for \(0\le t\le t_0\) are effective generalized CYs, klt for positive \(t\), \(d\ge0\) on all divisorials; take \(L_t\) affine, and \(t_0>0\) fixed rational.
Then eventually there is a uniform smaller \(t_1>0\) on which discriminants equal \(a'+t d'\) with \(a'\) of bounded denominator (i.e., the same \(p'\) works on all models).
There is also the following version with extra data:
\begin{enumerate}[label=\textup{(\roman*)}]
\item \(a+\lambda h+s c\) are effective generalized CYs over \(S\) for \(0<\lambda\le\lambda_0,\ 0<s\le\lambda\theta_0\), \(\lambda_0,\theta_0>0\) fixed, \(h\) of bounded denominator, \(c\ge0\);
\item on \(a=0\), \(d\le c\le C d\), \(-\eta_i c\le h\le Cc\) with fixed \(C\) and \(\eta_i\to0^+\).
\end{enumerate}

Then there are uniform \(\lambda_1,\theta_1>0\), with
\(\lambda_1\le\lambda_0\) and \(\theta_1\le\theta_0\), such that
eventually, on the sector
\(0<\lambda\le\lambda_1\), \(0<s\le\lambda\theta_1\), the data
are klt generically over \(S\), their discriminants are exactly
\(a'+\lambda h'+s c'\), and \(h'\) has bounded denominator.
Moreover, eventually there is \emph{no} divisorial \(T\) horizontal
over \(S\) with \(a(T)=0,\ h(T)<0\).
In particular for \(S\) a point exact nonnegativity on \(a=0\) follows even though only the approximate bound \(-\eta_i c\) was assumed.

In this statement the minima are lexicographic: first minimize \(a/e\),
then \(d/e\) on its minimum locus; for the extra data first minimize
\(a/e\), then \(h/e\), then \(c/e\).
Here \(e\ord_P=T|_{k(S)}\).
Thus \(a'\) itself is the discriminant in discrepancy convention even
if the endpoint is not klt generically.
The family \(a'+t d'\), \emph{including its endpoint at zero}, has the generalized effective log data (with nef b-part and total logarithmic divisor \(L_t\)).

\end{lemma}

The endpoint assertion follows from the positive-parameter assertion.
Indeed, for positive parameters use the generalized canonical bundle formula.
The moduli b-parts chosen using \(L_t\) are affine on \emph{all} models by the affine discrepancies; thus the limit b-divisor at zero is a difference of \(\mathbb Q\)-Cartier b-divisors by extrapolation, and nef by passage to the limit on a common determination.
The same applies on any intermediate base using pullback of \(L_t\).

We give details of the geometry and proofs of these lemmas.
Boundedness of denominators here and later does not require that the b-parts of the input have bounded Cartier index.

\subsection{Extraction and the Mori step}\label{subsec:A-mori}
We can pass to a small \(\mathbb Q\)-factorialization on which the
families still have effective traces.
Given an effective generalized klt CY \(A\) over \(S\) and Fano type, we can extract all exceptional divisors with \(A<\delta_*\) (\(0<\delta_*<1\)) to a \(\mathbb Q\)-factorial Fano type model over \(S\), extracting no others (or just select any subset of this list).
There are finitely many: on a log resolution where the nef part descends any further ones must be toroidal for the snc support (the integral reduced-log discrepancy otherwise contributes at least 1: all lc places of a reduced snc pair are monomial at the strata), and all their weights are bounded using positivity on the finitely many components.
For extraction one can pass to an ordinary big klt CY over \(S\) having the selected divisors still strictly positive in log-pulled boundary.
For clarity, mix the given data with small weight of auxiliary crepant sub-data with ample b-part determined on that resolution, horizontal and vertical trace on \(X\) effective and containing a general relatively ample contribution; these auxiliary data with total \(\sim_{\mathbb Q}0/S\) exist by taking the ample b-part sufficiently small and using a klt log Fano pair on \(X/S\).
More explicitly subtract the small pushforward of the ample divisor and a small ample divisor on \(X\) from the log Fano ample class (twist by a pullback if needed).
Sub-klt persists on a common resolution after mixing, as do effectivity on \(X\) and positive coefficients at the selected divisors.
Replace the resulting ample moduli contribution on the high model by general representatives (over \(S\), twisting by pullback), with small coefficients.
Apply klt extraction; the log transform is effective and still big over \(S\) by the general ample contribution on \(X\) avoiding the finitely many centres.
Big klt CY implies Fano type.

Other CY families retain effective traces on the extracted model
provided their discrepancies at the extracted divisors are at most
\(1\).  Each application below verifies this condition.
When the whole list is retained, the resulting \emph{ordinary variety}
is \(\delta_*\)-lc.  For an exceptional divisor on a higher
determination, effectivity of the trace and the negativity lemma for
the difference between the pullback of the nef-part trace and its nef
determination give \(A_X(T)\ge A(T)\).  Every divisor for which
\(A(T)<\delta_*\) is already on the extracted model; its ordinary
discrepancy there is \(1\).  The remaining exceptional divisors have
\(A_X(T)\ge A(T)\ge\delta_*\).
A \(K\)-MMP over \(S\), if \(\dim X>\dim S\), now ends at a \(\delta_*\)-lc Fano Mori step \(g:X_1\to W\) over \(S\) with \(\dim W<\dim X\), preserving the CY discrepancies/effectivity we arranged and Fano type.
Here \(-K_X\) is big over \(S\), and the ordinary discrepancies stay bounded below by monotonicity in the \(K\)-MMP; the geometric generic fibre \(F\) of \(g\) (we also write \(F\) for the fibre over \(k(W)\)) is a bounded Fano by BAB (generic discrepancies via horizontal divisors and separable constant extensions). \(W\) is Fano type over \(S\) (image of Fano type).
If instead the contraction to \(S\) is birational no further step is
needed: the function fields agree and the discriminant is the original
discrepancy function.  This is the terminal case in each induction.

\subsection{Tailored complements and connected minimizing places}

\begin{lemma}[Tools at an lc place]\label{lem:A-lc-place-tools}
Let \(g:X_1\to W\) be a Mori step as above, and let \(A_*\) be an
effective generalized CY discrepancy on \(X_1\) over \(S\).
\begin{enumerate}[label=\textup{(\roman*)}]
\item Suppose that \(A_*\) is lc over the generic point of \(W\),
and that an effective generalized klt CY over \(S\) is available for
perturbation.  Any prescribed horizontal divisorial lc place remains
an lc place of an ordinary boundary \(\Omega\), effective on the
generic fibre \(F\), with \(n(K_F+\Omega_F)\sim0\), for a bounded
positive integer \(n\).
\item For \(P\) on a birational model of \(W\), suppose \(A_*\) is klt generically and \(T\) computes the discriminant \(b_*(P)\).
There is an ordinary \(n\)-CY over the DVR, lc in its b-discrepancies
and effective horizontally on \(F\), with \(n\) bounded, keeping
\(T\) as an lc place.  Here \(n\)-CY means that the ordinary total
log canonical divisor, multiplied by \(n\), is linearly trivial
over the trait.
Moreover for the adjusted discrepancy \(A_*-b_*(P)e\), any two minimizing divisorials can be joined by a chain/interpolation of lc places in the following sense: on a log resolution over the generic DVR their strict divisors (or extractions of them) belong to a connected snc floor consisting of vertical coefficient-one divisors with monomial cones at the intersections.
One can include any additional finite discrepancy data in the resolution.
In particular normalized evaluations \(a/e\) of such other data interpolate continuously along the chain, using rational weights to approximate intermediate values by divisorial evaluations.
\end{enumerate}

\end{lemma}

\begin{proof}
Here everything over a generic DVR can be carried out by shrinking about a generic point of a prime in a smooth model over \(k\).
Only assertions over the generic fibre and that DVR are needed.
\emph{A model retaining the selected fibre components.}
Write $R$ for the trait ring.  Fix a nonempty finite collection of
divisors computing $b_*(P)$ and choose a log resolution $Z$ over
$\Spec R$ on which they occur.
For the adjusted generalized data $A_*-b_*(P)e$, they have boundary
coefficient one.  We first construct an ordinary klt pair on $Z$ whose
minimal model contracts every other special-fibre component and every
horizontal divisor exceptional over $F$.

Subtract a sufficiently small positive multiple of the fibre from the
adjusted boundary.  The resulting generalized subpair is sub-klt, and
its coefficients on the selected components remain positive.  Mix it
with a sufficiently small amount of auxiliary sub-CY data having an
ample b-part on $Z$ and effective nonexceptional horizontal trace.
Such data exist because $F$ is Fano: a small pushed ample contribution
can be subtracted from $-K_F$, and vertical terms can be adjusted freely
over the trait.  The pushforward is rationally Cartier on the
$\Q$-factorial total space $X_1$.
Replace the mixed ample b-part by general relatively free
representatives.  Make these choices on a common resolution and retain
a general big contribution avoiding the components whose coefficients
will be changed.

Increase the coefficients of all horizontal divisors exceptional over
$F$ and all special-fibre components outside the selected collection,
strictly and to values in $[0,1)$.  We obtain an effective ordinary klt
boundary $\Delta$, big over the trait, and
\[
                    K_Z+\Delta\sim_{\Q,R}E,
\]
where $E\ge0$ is supported precisely on the divisors just increased.
In particular at least one special-fibre component is absent from
$\Supp E$.

Here the relative fixed part can be computed directly.  For divisible
$m$, a rational section of $\mathcal O_Z(mE)$ restricts on the generic
fibre to a rational function with poles only along divisors exceptional
over the normal projective variety $F$.  It descends to a regular
function on $F$, and is therefore a constant in the trait field.
The section condition on a retained special-fibre component, whose
coefficient in $E$ is zero, forces that constant to have nonnegative
DVR order.  Conversely every element of $R$ gives such a section.
Thus the relative section module is $R$, and the relative fixed part
of $|mE|$ is exactly $mE$.

The good minimal model theorem for klt pairs with big boundary now
applies, after spreading over an open neighbourhood of the generic
point of $P$.  Every curve in a negative extremal ray lies in $\Supp E$, so
no divisor outside this support is contracted.  The section ring is
preserved by the MMP.  On its semiample output the surviving trace of
$E$ must therefore vanish, since otherwise it would remain a nonzero
fixed part.  Consequently every component of $E$ is contracted.
Call the output $Y$.  Its boundary remains big and klt and its log
canonical divisor is rationally trivial over the trait, so $Y$ is of
Fano type there.  Its generic fibre agrees with $F$ in codimension one.
The original adjusted generalized data are crepant through these
birational operations.  Their trace on $Y$ is effective: the only
special-fibre components are the retained ones, with coefficient one,
and the surviving horizontal coefficients are the effective
nonexceptional coefficients over $F$.

\emph{Connectedness of the minimizing locus.}
Let $\pi:\widetilde Y\to Y$ be a log resolution on which the nef part
of the original adjusted data descends.  Write $B^{\log}$ for its
boundary and $H=(B^{\log})^{=1}$.  The negative boundary is
$\pi$-exceptional because the trace on $Y$ is effective.  Moreover
$-(K_{\widetilde Y}+B^{\log})$ is nef over $Y$, and is big over $Y$
because $\pi$ is birational.  The birational connectedness lemma
therefore makes $H$ connected over each fibre of $\pi$.

One can see the relevant structure-sheaf assertion directly.  Put
$G=\lceil-(B^{\log})^{<1}\rceil$, an effective exceptional divisor.
Relative Kawamata--Viehweg vanishing for
$\mathcal O_{\widetilde Y}(G-H)$ gives a surjection
\[
               \mathcal O_Y\longrightarrow\pi_*\mathcal O_H(G).
\]
The image factors through $\pi_*\mathcal O_H$, which injects into
$\pi_*\mathcal O_H(G)$.  Hence $\pi_*\mathcal O_H$ is the structure
sheaf of the reduced image.  This image is the entire reduced special
fibre of $Y$, since all its components have coefficient one.  The
special fibre is connected by Stein factorization.  Thus $H$ is
connected, and the prescribed divisors lie in one connected snc
coefficient-one locus.  Generic klt ensures that every component of
this locus is vertical.  Monomial valuations along its intersections
join its divisorial orders continuously; the uniformizer has positive
value throughout.  Evaluations of any additional discrepancy data are
linear on these cones after taking a common resolution, so their
ratios to $e$ interpolate continuously.  Rational weights recover the
required divisorial interpolants.

\emph{An ordinary complement preserving the chosen place.}
For a tailored ordinary \(\Omega\), use the model above (or horizontally, extract the horizontal place using a CY slightly perturbed towards klt as funding).
Write \(T_0\) for the resulting prime, which has coefficient one.
Run a \(-(K+T_0)\)-MMP; this divisor is pseudo-effective over the trait (respectively generically) by the effective generalized CY containing \(T_0\).
On every model the pair with strict \(T_0\) is lc by comparison with the generalized data and b-nef negativity.
At the nef model take a bounded ordinary lc complement (coefficients only 0 and 1, Fano type; in the relative case use Birkar's complement over a closed point of the relevant prime/open, after appropriate shrinkings as above, containing the generic trait/field in the resulting neighbourhood).
Indeed over the trait/field the final nef divisor is semiample by Fano type; generation of a multiple spreads out, as does a Fano type witness (relative ampleness and klt) and lc of the indicated pair via resolution, shrinking by projectivity to ensure the conditions for complements over an open.
Its crepant pullback boundary on the model with \(T_0\) still dominates \(T_0\): it adds the pull of the effective complement addition and the effective MMP discrepancy of \(-(K+T_0)\).
In particular contraction of \(T_0\) in a negative step would even contradict lc.
Push generically to \(F\), preserving effective boundary there.
All complement indices use only the absolute dimension bound.
This proves both assertions.

\end{proof}

\subsection{Uniform charts and discrepancy estimates}

\begin{proposition}[Uniform charts]\label{prop:A-uniform-charts}
At a chosen vertical place in Lemma~\ref{lem:A-lc-place-tools}\textup{(ii)},
after a bounded degree extension of \(k(W)\) and a chosen extension of
the DVR and of the place, there is a toroidal cone representing the
place, with nonnegative coordinates
\begin{equation}\label{eq:A-contact-cone}
       (x,e_*),\qquad M x=r e_* .
\end{equation}
Here \(M\) ranges over a finite list of integral matrices of bounded sizes, \(r\) is integral (possibly huge), \(e_*\) is value of the new uniformizer, and original uniformizer value is \(\rho e_*\) with positive bounded \(\rho\).
Nonzero integral solutions on the cone give divisorial valuations (not necessarily primitive after restriction); evaluations of discrepancy data always refer to homogeneous values by restriction.
For horizontal places as in the first bullet one just has an orthant on a bounded generic resolution, without \(e_*\).
For \emph{any} effective generalized CY data \(A\) on \(X_1\) over \(S\), values on these cones are finite continuous homogeneous piecewise linear, and are Lipschitz in \(x\) with uniform constant \emph{at fixed} \(e_*>0\) (uniform homogeneous Lipschitz in the horizontal case).
Constants depend on the boundedness in this Mori step and the complement indices, \emph{not} on \(r\), vertical coefficients, or the nef part index or the chosen data \(A\).
In particular finitely many differences at fixed parameter separations have uniform Lipschitz bounds as well.

\end{proposition}

The proof separates the construction of the charts from the metric
estimate.  The first lemma keeps track of their integral lattices and
logarithmic coordinates; the second converts generic-fibre degree into
an angular bound.  We then replace the nef part at two chosen valuations
to apply that bound to arbitrary generalized CY data.

\begin{lemma}[Bounded toroidal charts]\label{lem:A-toroidal-charts}
For the prescribed place and tailored complement in
Proposition~\ref{prop:A-uniform-charts}, there are charts with its stated
contact equations and integral-valuation property.  The prescribed
place is monomial on one of these charts.  They have rational functions
\(z_1,\ldots,z_f\), \(f=\dim F\), with uniformly bounded pole degrees on
a bounded embedding of \(F\), such that:
\begin{enumerate}[label=\textup{(\roman*)}]
\item their logarithmic differentials form a relative log basis in the
vertical case and an absolute log basis over the generic field in the
horizontal case;
\item their orders are linear functions of the chart coordinates with
uniformly bounded slopes;
\item every rational monomial valuation in the vertical chart with
\(e_*>0\) restricts to a Gauss valuation on
\(K_1(z_1,\ldots,z_f)\), where \(K_1/k(W)\) is the bounded extension
used in constructing the chart.  Horizontally the same assertion holds
with the trivial valuation on the generic field.
\end{enumerate}
Here Gauss means that evaluating a polynomial in the \(z_j\) gives the
minimum of its coefficient-adjusted monomial weights, without
cancellation among terms of equal weight.
\end{lemma}

\begin{proof}
\emph{Bounded marked fibres and descent of an embedding.}
The pairs \((F,{\rm Supp}\,\Omega_F)\) are in bounded embedded families over the geometric field.
Indeed boundedness gives a very ample class \(\mathcal H\) of bounded degree and embedding dimension on each bounded \(F\).
One can use the complete system here (dimension still bounded by degree for the nondegenerate image).
To bound the anticanonical degree, let \(f=\dim F\) and take a
smooth complete-intersection curve \(C\) of \(f-1\) general members
of the very ample system \(\mathcal H\), avoiding the singular
locus.  Adjunction gives
\[
 -K_F\cdot\mathcal H^{f-1}
   =(f-1)\mathcal H^f+2-2g(C)
   \le (f-1)\mathcal H^f+2 .
\]
For \(f=1\), the same formula is the usual degree formula on the
smooth curve \(F\).
The nonzero coefficients of \(\Omega_F\) are at least \(1/n\),
so
\[
 \deg_{\mathcal H}\Supp\Omega_F
 \le n\deg_{\mathcal H}\Omega_F
 =-n K_F\cdot\mathcal H^{f-1}.
\]
The support therefore has bounded degree.
In the DVR case we can secure the embedding over a bounded degree extension of \(k(W)\).
The geometric Picard group is free abelian of bounded rank: it injects by normality into the Picard group of a smooth projective resolution, which is rationally connected (rational connectedness for klt Fano) and has bounded Betti numbers for a suitable resolution, by resolving the bounded families after stratification.
We use here the usual characteristic zero Picard consequences of smooth rational connectedness (in particular trivial connected Picard and no torsion).
The finite image of Galois on the discrete Picard lattice has bounded
order.  One way to see the uniformity is to reduce a finite subgroup of
\(\operatorname{GL}_r(\Z)\) modulo \(3\): the principal congruence
subgroup at \(3\) is torsion free, so this reduction is injective, and
\(\lvert\operatorname{GL}_r(\mathbf F_3)\rvert\) bounds its order.
After bounded extension the chosen \(\mathcal H\) is invariant; its complete linear system descends projectively as a projective-space form of bounded dimension (the ambiguity on section maps is scalar), which splits after another bounded extension (Severi--Brauer index bound).
This suffices.
Bounded-degree embedded varieties and their marked reduced divisors as
above are parametrized by finitely many Chow families
\cite[Definition~3.21, (3.21.1)--(3.21.3)]{KollarChow2017}.
Here we use effective cycles and their reduced supports, with the
universal cycles restricted to smooth characteristic-zero strata.

\emph{Toroidalization of the fixed families.}
Apply marked weak toroidalization
\cite[Theorem 1.1]{AbramovichDenefKaru} to these families.
Its source and base modifications can be taken projective and
birational; allowing a fixed alteration would give the same construction
after a bounded extension.
Here are the uniformity and chart details we use.
By compactification and noetherian stratification, for a dense open of each parameter stratum we can use a single proper toroidalized family over a modification (an alteration would also suffice) of its compactified base.
Over the dense open (which can always be shrunk, treating its complement by noetherian induction) the new fibres give birational models of the original varieties.
Here one works stratum by stratum starting with geometrically integral generic fibres, retaining these and birationality in the spread-out open (in case of base alteration it suffices to modify the base-changed family birationally).
The hypotheses on the fibres parametrized, such as geometric integrality and marked support of lower dimension, can be imposed by ordinary Hilbert stratification.
The compactification of a stratum serves to extend traits whose \emph{field} point is in that dense open; their closed points need not parametrize good pairs at all.
We include horizontal markings as follows: first resolve the generic fibre with its marked support over this stratum, spread out over an open, and mark the strict boundary and the exceptional locus; arrange their preimage in the toroidalization to be in the boundary.
Thus on the resulting generic-fibre models relevant to us any \emph{positive} crepant coefficient of \(\Omega_F\) is supported on the toroidal boundary: outside the old marking one had smooth with crepant subboundary \(\le0\).
We arrange the dense open to be interior on the toroidal base of the family.
We may use smooth strict toroidal models (also by toroidal resolution); all modifications here are of fixed finite type families.
Any base alteration only requires bounded degree field extension to lift the field point on the chosen dense open.
Thus after bounded degree extension of \(k(W)\), say to \(K_1\), and extension of the place, we have a lift of the whole new trait to this base by properness, with its generic point interior.
Write \(R\) for the new trait ring, \(t_*\) for its uniformizer, \(\rho\) the ramification index over the original DVR (bounded).
We always use log pullback over the finite extension.
One may view the trait as the localization at a prime divisor in a smooth model over \(k\), working close to the generic point of that divisor (and spreading the map to such a neighbourhood, shrinking to pull the base boundary only to the closed-point boundary there).
In particular the divisorial log structure of the trait comes from a log-smooth \(k\)-pair locally.

\emph{Full character rank and the log chart.}
We verify log smoothness at the strata, including the case in which the
map on boundary exponent groups has a kernel.
In a completed toroidal description use the full character lattices of
the toric models.  These include invertible torus characters as well as
the characters cutting out the boundary.  Dominance gives an injective
map of full character lattices after tensoring with \(\Q\).  In their
logarithmic differential bases its matrix has constant integral
entries.  It therefore has full target rank at the closed stratum as
well as at the generic point.  Any lattice torsion is tame in
characteristic zero.  The unit characters are essential here: two
target boundary characters can have the same boundary exponents on the
source while their quotient pulls back to a nonconstant unit, whose
differential supplies the missing direction.

To detail the algebraic chart check, at closed points in these smooth strict models write the target boundary equations \(y_j,\ 1\le j\le s\), as pulling to \(u_j x^{\alpha_j}\) with \(x_1,\ldots,x_m\) source boundary parameters and \(u_j\) units.
The map of log differential fibres has full target rank (as above in the completed toroidal models, where boundary equations match up to units).
For \(\alpha:\mathbb Z^s\to\mathbb Z^m\) the exponent map, let
\(K=\ker(\alpha)\).
The quotient \(\Z^s/K\) is the free group \(\operatorname{im}(\alpha)\),
so \(K\) is saturated and admits an integral projection
\(p:\Z^s\to K\).
The unit homomorphism \(u\) from \(\mathbb Z^s\) factors as its kernel-direction contribution via this projection, times a homomorphism on \({\rm im}(\alpha)\).
The latter extends to \(\mathbb Z^m\) after adjoining unit roots \'etale locally, so can be absorbed by rescaling the \(x_\ell\) by units.
We now have the chart \(\mathbb N^s\to \mathbb N^m\oplus K\) using \(\alpha\) and the projection, with the free group summand mapping on the source by unit functions \(u|_K\).
This is injective on groupifications with tame cokernel torsion.
The induced map to the product of the target with the source toric chart scheme over \(\mathbb A^s\) is smooth near the point: besides the rescaled \(x_\ell\), one needs independent ordinary differentials for those kernel units and the pulled target transverse coordinates along the source stratum.
Exactly these are independent there by full log differential rank (their boundary-weight entries are zero).
This proves the smoothness required in Kato's chart criterion
\cite[Theorem~3.5 and Remark~3.6]{KatoLog}; see also the statement and
proof in \cite[Theorem~4.1, Remark~4.2 and Section~6]{FumiharuKatoLog}.
In particular the required rank at a closed stratum follows from the
constant full-character matrix; it does not require specializing an
arbitrary ordinary differential minor.

\emph{The trait and its integral contact cone.}
Take the fine saturated log fibre product with the trait equipped with
its closed-point divisorial log structure.  Log smoothness is preserved
by fs base change and composition \cite[Section~1.1]{IllusieNakayamaTsuji}.
Spread the trait to the log-smooth $k$-pair described above and form the
same fs fibre product there.  Its source is log smooth over $k$ with
the trivial log structure.  It therefore admits smooth local charts
over affine toric varieties
\cite[Theorem~4.7 and Section~7]{FumiharuKatoLog}.
In these charts the monomial stratum quotient is regular, and its
dimension plus the characteristic-monoid rank is the ambient dimension.
Thus the spread-out source is log regular; localization gives this
property for the fs fibre product over the trait.
Its main component is proper over the trait and gives the desired generic model.
Indeed the map from the fs fibre product to the ordinary one is proper here (in charts, integralization and finite toric saturation); the trivial-log open is dense in a log regular model, and maps to the generic point of the trait.
At a stratum with boundary generators from the total family and from the trait, the cone of nonnegative orders is given precisely by \(Mx=r e_*,\ x,e_*\ge0\).
Here rows come from the toroidal base boundary equations, pulled to total boundary monomials times units with multiplicities \(M\), and to the trait with contacts \(r\); thus \(M\), unlike \(r\), comes from fixed chart types.
These are the matching conditions from the pushout chart (the total family and its base before this log base change can have smooth strict toroidal structures, hence boundary-coordinate charts).
After spreading near the trait the base has also ordinary residue/transverse coordinates, with the prime divisor as its single boundary near the generic point in use; only its boundary weight thus enters these equations.
Units have weight zero; at the point used only boundary generators vanishing at its images in the factors are counted.
No further independent zero conditions on weights arise from units of the saturated monoid becoming units at the point: a positive multiple of such an element comes from the original generators and if it is a unit cannot involve any generator vanishing at the point.
An integral match defines an integer-valued homomorphism on the
groupification of the pushout monoid.  Saturation takes place in that
same groupification, so this homomorphism remains integer-valued on
the saturated monoid.
We may work in strict \'etale local charts at a geometric stratum and then restrict valuations.
Thus nonzero rational weights give monomial divisorial evaluations (and the zero vector the trivial evaluation); integral ones restrict to integral multiples of primitive divisorials when nontrivial.
Indeed in the split toroidal chart one is taking the dual cone of the characteristic monoid at the stratum.
A nonzero rational ray even on a face can be introduced in a toroidal subdivision (over the corresponding generization stratum for a face); the primitive vector for its toric lattice then gives a boundary-divisor order.
Matched integral weights as just described are integral in this lattice.
For \(e_*>0\) the restriction to \(K_1\) is \(e_*{\rm ord}_{R}\).
This description is equivalently obtained by fs pulling back the log-smooth monomial charts and subdividing toroidally.
Use the family of monomial valuations for one chosen lift of the stratum in a split local chart (cones/subdivisions there, then restrict to the original function field); integral points give integral orders since they already do so in such a chart.
No degree bound on further strict \'etale local charts is needed here:
restriction of an integer-valued valuation still has value group a
subgroup of \(\Z\), hence is an integral multiple of a primitive
divisorial order.  It is never divided by the chart degree.
Only the initial field extension contributes the separately bounded
ramification index \(\rho\).  The rational functions used in the metric
estimate below are taken on the fixed total family itself.

All special components are log boundary on this model.
The subboundary of the crepantly pulled ordinary \(\Omega\) is bounded above by the reduced log boundary (coefficient bound by lc, horizontal positive support marked as above).
Thus its chosen lc place is a place also of the reduced-log pair, since the difference is effective \(\mathbb Q\)-Cartier.
Such lc places are monomial on the log regular model: pass to a toroidal snc subdivision and use the lc place description for an snc pair.
We take a cone as above containing the extended chosen place.
For a horizontal place as in the first bullet one just takes a bounded stratified log resolution of the generic pairs \((F,{\rm Supp}\,\Omega_F)\); same reasoning yields an ordinary snc monomial cone (over the generic field, or its extension).

\emph{Logarithmic coordinates and Gauss restriction.}
For the later metric estimate, choose rational log coordinates
\(z_1,\ldots,z_f\), \(f=\dim F\), on charts of the \emph{fixed}
toroidalized families, such that \(d\log z_j\) form a basis of relative
log differentials near the stratum in question.
Indeed the latter are locally free by log smoothness, and one can choose among boundary parameters and invertible functions giving generators from the smooth total chart.
Finitely many choices of genuine rational functions on opens of smooth strict charts suffice (Cartier local boundary equations and local units generate logarithmic differentials).
Their weights are bounded linear forms in the cone coordinates (boundary parameters or units), and the basis property persists under fs base change and toroidal resolution.
Moreover the \(z_j\) have bounded pole degrees on the embedded \(F\)'s for these charts: they are chosen on fixed families birational to the embedded ones, so on a sufficiently shrunk dense open of each stratum they give fixed bounded-degree rational maps on the embedded fibres.
This shrinking can be built into the stratification (express each function by fractions of bounded-degree embedded coordinates on the generic fibre of the family and spread).
Use the analogous absolute log bases in the horizontal generic case.

A rational toroidal valuation here with \(e_*>0\) restricts to a Gauss valuation on \(K_1(z_1,\ldots,z_f)\).
To check, rescale to a primitive divisor on a toroidal subdivision.
Write \(\nu_j=v(z_j)\in\mathbb Z\), \(e_*=v(t_*)>0\).
At this divisor the residues of \(z_j^{e_*}t_*^{-\nu_j}\) have independent log differentials.
Indeed relative log cotangent there, modulo the divisor, is identified with residue differentials over the base residue field: \(t_*=\pi^{e_*}\cdot(\text{unit})\), for uniformizer \(\pi\), eliminates \(d\log\pi\), while base unit differentials restrict as usual.
This also follows from the exact sequences for log differentials at the generic smooth divisors (and separability).
The indicated unit log differentials are \(e_*\) times the chosen basis in relative log cotangent.
The residues are thus algebraically independent.
In particular a lattice basis of the weight ties (monomials \(z^\gamma\) with \(\nu\cdot\gamma\) a multiple of \(e_*\), adjusted by that power of \(t_*\)) has algebraically independent residues too, so no cancellation occurs at equal coefficient-adjusted weights.
Horizontally one keeps \(d\log\pi\), without a base uniformizer; for a nontrivial monomial valuation of integral rank one weights, a basis of the kernel characters of the weight gives independent residues by the same log basis calculation.
This again gives Gauss, now for the trivial base valuation.

\end{proof}

\begin{lemma}[Degrees control angular variation]\label{lem:A-function-slopes}
On a chart of Lemma~\ref{lem:A-toroidal-charts}, the order of a nonzero rational
function is continuous, homogeneous and finite piecewise linear,
including on the faces.  If its pole divisor on the embedded generic
fibre has degree at most \(m\), where \(m\ge1\), its Lipschitz constant
in \(x\) at fixed \(e_*>0\) is at most \(Cm\).  Horizontally this bound
holds on the entire cone.  The constant \(C\) depends only on the fixed
families and chart choices.  A function with pole degree less than one
is constant on the generic fibre and has zero angular variation.
\end{lemma}

\begin{proof}
\emph{Piecewise linearity and faces.}
On these cones rational functions have continuous homogeneous finite piecewise linear values, including faces.
For example subdivide a split chart toroidally to snc (subfaces using the corresponding generization strata in the subdivision) and use the monomial order formula at the generic stratum (after splitting a chart if necessary), by expanding at the snc parameters: minima on power series supports are finite polyhedral by monomial noetherianity.
The same minima using weights zero off a subcollection of parameters give the values on its face, i.e. at the more general stratum.
Indeed the monomial ideals in the retained parameters test order there; membership is still tested in the power series ring localized along the prime of those parameters (faithful flatness and the monomial formula; a positive-threshold monomial ideal for the retained parameters only is primary to their prime, as seen already in completion).
This gives continuity and compatibility across subdivisions.
In a local split toroidal chart the subdivisions are toric subdivisions at that stratum; one is using the corresponding monomial order family throughout (weight ideals of boundary monomials), so that subfaces common to two cones use the same such orders.

\emph{A degree bound for angular slopes.}
Suppose a nonzero rational function \(g\) on \(F\) over the extended field has pole degree \(\le m,\ m\ge1\), in the bounded embedding.
It satisfies a polynomial equation over the \(z\)'s of bounded degree in \(g\) and degree \(O(m)\) in the \(z\)'s.
Use the equation of the image hypersurface of \((z,g)\) in a product of projective lines.
Its multidegrees are bounded by intersecting general level divisors of all but one coordinate on \(F\) on the locus of definition: individual divisor degrees are bounded by their pole degrees, so projective B\'ezout bounds isolated general solutions by a constant times the product bound (characteristic zero).
Equivalently intersect the product of those level-divisor supports with diagonal conditions of bounded degree in a product of projective spaces, counting isolated points (the product variety has Segre degree bounded by a fixed multiple of the product of individual degrees).
Write this equation as
\[
 \sum_{j=0}^q c_j(z)g^j=0,
 \qquad q\le C,\qquad \deg_z c_j\le Cm .
\]
By the Gauss property, \(v(c_j)\) is a minimum of affine functions whose
\(x\)-slopes have norm at most \(Cm\).  Their constant terms can involve
arbitrary orders of coefficients in the trait field, but at fixed
\(e_*\) these terms do not vary.
At least two nonzero terms of the equation attain its least valuation.
Consequently \(v(g)\) is one of the finitely many quantities
\[
 \frac{v(c_i)-v(c_j)}{j-i}\qquad(i\ne j).
\]
Each is \(Cm\)-Lipschitz after enlarging the fixed constant.
Along a segment the continuous, finite piecewise linear function
\(v(g)\) selects from this finite list; subdividing at its finitely many
switches gives the same Lipschitz bound on the entire segment.
Rational weights are dense, so the estimate holds also for real
weights.  Horizontally the argument applies on the entire cone.
If the pole degree itself is \(<1\) the function is constant along the generic fibre.

\end{proof}

\begin{proof}[Proof of Proposition~\ref{prop:A-uniform-charts}]
Lemma~\ref{lem:A-toroidal-charts} supplies the chart, the prescribed
monomial place and the integral orders.  We first extend the
piecewise-linearity assertion of Lemma~\ref{lem:A-function-slopes} from
rational functions to the discrepancy data.
Evaluations of \(\mathbb Q\)-Cartier b-divisors are also finite piecewise linear: locally over an affine neighbourhood of the stratum write a determining Cartier multiple on a projective model as a difference of relatively generated divisors, then it uses minima on finite sets of rational functions.
Generalized discrepancies differ on this cone from the reduced-log discrepancy (identically zero) by such Cartier b-evaluations.
Piecewise-linearity here need not have bounded combinatorial complexity.


\emph{Two-point replacement of the nef part.}
Fix two rational valuations on one cone with the same \(e_*\), and
choose a common high model on which their centres and the determination
of the nef part are visible.  On that model let \(L\) be an ample
divisor, relative to the trait, and choose a rational
\(\varepsilon>0\).  After twisting by a pullback from the base, take
general effective representatives
\[
 E_+\sim_{\Q}\mathbf M+\varepsilon L,\qquad
 E_-\sim_{\Q}\varepsilon L
\]
which contain neither centre.  The ample perturbation is selected
\emph{after} the high model: its coefficient can therefore be made so
small that its pushed horizontal degree is arbitrarily small.
Add \(E_+-E_-\) to the crepant generalized boundary and remove the nef
part.  Since both representative divisors have order zero at both
chosen valuations, the resulting ordinary divisor data have exactly the
prescribed discrepancies at these two valuations.

Let \(\Theta\) be its horizontal trace on the bounded fibre \(F\).
Write degrees with respect to the bounded very ample class.
The original trace \(B_F\) is effective, and the trace of the nef part
has nonnegative degree.  Relative triviality gives
\[
 \deg B_F+\deg\mathbf M_F=-\deg K_F .
\]
Negative exceptional crepant coefficients on the high model disappear
when pushed to \(F\).  Hence, if \(L_F\) denotes the pushed ample
contribution, then
\[
 \deg\Theta^+\le -\deg K_F+\varepsilon\deg L_F,
 \qquad
 \deg\Theta^-\le \varepsilon\deg L_F .
\]
Both degrees are uniformly bounded.  This estimate uses no bound on
the height of the determination or the denominator of \(\mathbf M\).
For divisible \(l\), relative rational triviality supplies a rational
\(l\)-pluriform \(\sigma\) for these ordinary data.  Its divisor on the
generic fibre is \(-l\Theta\); the possible pullback term from the trait
has constant value on a slice with fixed \(e_*\).
Divide by the \(l\)-power of the log form \(\psi\) from wedging our relative \(d\log z_j\)'s with a base log volume near this DVR.
At either evaluation the discrepancy is \(v(\sigma/\psi^l)/l\) up to the base term constant at fixed \(e_*\), since \(\psi\) is a log generator at the toroidal divisors.
The generic pole degree of this ratio is \(O(l)\): the log volume of the \(z_j\) has controlled poles (on the regular codimension one locus of \(F\)), hence controlled zeros by the canonical degree (nonpositive here by log CY effectivity horizontally), and the horizontal boundary degree was just bounded.
Lemma~\ref{lem:A-function-slopes} applied to \(\sigma/\psi^l\), followed by division
by \(l\), proves the uniform metric assertion.  The integer \(l\)
may be arbitrarily large.  Continuity extends the estimate from the
two rational endpoint valuations to any two real valuations on the
slice.  Thus neither the nef-part index nor the vertical contact
enters the constant.
Horizontal traces and these canonical-degree computations on the generic fibre can use forms divided by a separating base volume (the absolute log order at a horizontal divisor is the fibre log order, by taking independent base residues).
Horizontally use exactly the same replacement and log volume over the generic field.

\end{proof}

\begin{corollary}[Horizontal variants]\label{cor:A-horizontal-variants}
The chart and metric conclusions have the following two variants.
\begin{enumerate}[label=\textup{(\roman*)}]
\item
The same tailorable bounded complement and homogeneous chart argument works if we are working over an algebraically closed generic field, with the relevant effective CYs on a \(\mathbb Q\)-factorial model of Fano type, and have boundedness only up to isomorphism in codimension one of these models.
\item
A Cartier rational b-order function \(N\) with effective trace of degree
\(O(\xi)\) on the bounded model, b-linearly \(\mathbb Q\)-equivalent
to a difference of b-nef parts both of trace degrees \(O(\xi)\), has
horizontal Lipschitz constant \(O(\xi)\), including when \(\xi\to0^+\).
\end{enumerate}
\end{corollary}

\begin{proof}
For \textup{(i)}, tailor the complement on the Fano type model using
the klt perturbation in Lemma~\ref{lem:A-lc-place-tools}, and then push
the complement and the effective CY data crepantly to the bounded
normal small model.  Their total divisors are \(\Q\)-linearly trivial.
This supplies the bounded embedded marked family and all the degree
estimates in Proposition~\ref{prop:A-uniform-charts}; the
anticanonical degree is bounded by the same curve slicing argument.

For \textup{(ii)}, use the two-point construction for each of the two
b-nef parts, taking the pushed ample errors to be \(o(\xi)\).
As above replace those parts by general effective representatives after the same arbitrarily tiny ample addition on a high determination, avoiding both centres.
The residual rational principal term for \(N\) now uses functions with pole degree divided by denominator \(O(\xi)\) on \(F\); apply Lemma~\ref{lem:A-function-slopes} (using constancy if pole degree \(<1\)).
\end{proof}

\subsection{Integral approximation on the cones}

The charts now give a bound on how discrepancy changes with the angular
coordinates, even when the trait contact is large.  We next approximate
real chart points by integral ones.  Integrality makes the resulting
orders positive integral multiples of prime orders, so the small-prime
hypotheses in the comparison lemmas can be used without changing their
normalization.

The following approximation does not require uniform smoothness over the
DVR.  All norms are fixed Euclidean norms; the matrices range over a
fixed finite list, so their kernel lattices have uniformly equivalent
coordinate norms.

\begin{lemma}[Integral approximation]\label{lem:A-lattice-approximation}
\begin{enumerate}[label=\textup{(\roman*)}]
\item
Let \(u_i\) lie in a fixed nonnegative orthant and
\(\lvert u_i\rvert\to\infty\).  After passing to a subsequence there
are \(s_i\to0^+\), nonzero nonnegative integral vectors \(n_i\), and
a fixed, possibly subsequence-dependent, \(C_*\) such that
\begin{equation}\label{eq:A-horizontal-approximation}
 \lvert n_i-s_i u_i\rvert\le C_*s_i,
 \qquad s_i\lvert u_i\rvert\longrightarrow\infty .
\end{equation}
\item
Let \(M_i x_i=r_i\), where \(x_i\ge0\), \(r_i\) is integral and
\(M_i\) belongs to a fixed finite list of integral matrices.
If \(\alpha_i>0\) tends to zero, then after passing to a subsequence
there are positive integers \(E_i\), nonnegative integral \(n_i\), and
a fixed \(C_*\) such that
\begin{equation}\label{eq:A-vertical-approximation}
 M_i n_i=r_i E_i,\qquad
 s_i=E_i\alpha_i\longrightarrow0^+,\qquad
 \lvert n_i-E_i x_i\rvert\le C_*s_i .
\end{equation}
\item
For every sufficiently small fixed \(\gamma>0\), the same matching
equations admit \(n_i\ge0\) integral and \(1\le E_i\le J(\gamma)\)
with
\begin{equation}\label{eq:A-bounded-dirichlet}
 \lvert n_i-E_i x_i\rvert\le\gamma ,
\end{equation}
where \(J(\gamma)\) is uniform in \(i\) and in the contact vectors
\(r_i\).
\end{enumerate}
\end{lemma}

\begin{proof}
For \textup{(i)}, pass diagonally so that for every integral character
\(\gamma\), the numbers \(\langle\gamma,u_i\rangle\) either converge
in \(\R\) or diverge in absolute value.  The characters with bounded
values form a subgroup \(\Gamma\).  Their limits are additive and
extend to evaluation on a vector \(u_*\).
For a fixed interval \(I=[a,b]\subset(0,\infty)\), the probability
distribution of \(s(u_i-u_*)\) modulo the lattice, with \(s\) uniform
on \(I\), tends to Haar measure on \(\Gamma^\perp\).
Indeed its Fourier coefficient at a character in \(\Gamma\) tends to
one; for a character outside \(\Gamma\), integration of the
exponential bounds the coefficient by
\[
 \frac{2}{2\pi(b-a)
       \lvert\langle\gamma,u_i-u_*\rangle\rvert}\longrightarrow0.
\]
Trigonometric approximation proves weak convergence.
Every neighborhood of zero in \(\Gamma^\perp\) has positive Haar
measure.  Choose intervals tending to zero sufficiently slowly and
neighborhoods of radius \(o(s_i)\), while also requiring
\(s_i\lvert u_i\rvert\to\infty\).
A diagonal choice gives an integral \(n_i\) with
\(\lvert n_i-s_i(u_i-u_*)\rvert=o(s_i)\), proving
\eqref{eq:A-horizontal-approximation}.
The error is eventually less than one in each coordinate.
Since \(s_i u_i\ge0\), an integral coordinate of \(n_i\) cannot be
negative.  Its norm tends to infinity, so \(n_i\ne0\).

For \textup{(ii)}, pass to a subsequence with \(M_i=M\).
Elementary integer linear algebra gives a rational \(x_{0,i}\) with
\(Mx_{0,i}=r_i\) and \(l x_{0,i}\) integral for a positive integer
\(l\) depending only on \(M\).
For example, solve using a fixed maximal nonzero minor and take \(l\)
divisible by its determinant.
Put \(H_0=\ker M\), \(\Lambda=H_0\cap\Z^m\), and
\(U_i=x_i-x_{0,i}\).
On the dual lattice of \(\Lambda\), pass diagonally so that every
normalized distance
\[
 \frac{\operatorname{dist}
   (\langle\gamma,lU_i\rangle,\Z)}{l\alpha_i}
\]
converges to a finite value, with a signed residual, or tends to
infinity.  The characters in the first case form a subgroup
\(\Gamma\).  Their signed residual limits are additive; choose
\(u_*\in H_0\) whose evaluations give these limits.

Sample positive integers \(j\) for which
\(s=jl\alpha_i\in[a,b]\), with \([a,b]\subset(0,\infty)\) fixed.
The distributions of \(jl(U_i-\alpha_i u_*)\) in
\(H_0/\Lambda\) tend to Haar measure on \(\Gamma^\perp\), which may
be disconnected.  For \(\gamma\in\Gamma\) the frequency, reduced
modulo \(\Z\), is \(o(\alpha_i)\), so its Fourier coefficient tends
to one.  Otherwise the geometric-series bound is at most
\[
 \frac{C\alpha_i}
 {\operatorname{dist}
   (\langle\gamma,l(U_i-\alpha_i u_*)\rangle,\Z)}
 \longrightarrow0 .
\]
Use neighborhoods of zero and a slow diagonal choice as before.
There are \(j_i>0\) and \(k_i\in\Lambda\) such that, with
\(s_i=j_i l\alpha_i\to0\),
\[
 \lvert k_i-j_i l(U_i-\alpha_i u_*)\rvert=o(s_i).
\]
Set \(E_i=l j_i\) and \(n_i=E_i x_{0,i}+k_i\).
These are integral matches and
\(\lvert n_i-E_i x_i\rvert\le C_*s_i\).
The same coordinate argument proves \(n_i\ge0\).

For \textup{(iii)}, apply simultaneous Dirichlet approximation to
\(lU_i\) on the fixed kernel torus.  It supplies a positive
\(j_i\le J_0(\gamma)\) with
\(\operatorname{dist}(j_i lU_i,\Lambda)\le\gamma\).
Use the same definition of \(n_i\) and \(E_i=l j_i\).
Matching is exact, positivity follows when \(\gamma<1\), and the
bound on \(E_i\) is independent of \(r_i\).
\end{proof}

In a horizontal discrepancy application take \(u_i=w_i/a_0(w_i)\).
In a vertical application of \eqref{eq:A-contact-cone}, take
\(x_i=w_i/e_{*,i}\) and
\(\alpha_i=a_0(w_i,e_{*,i})/e_{*,i}\).
The constant \(C_*\) in the vanishing-error conclusions is allowed to
depend on the selected subsequence; the bound \(J(\gamma)\) in
\eqref{eq:A-bounded-dirichlet} is uniform.

\subsection{Proofs of the comparison statements}

We apply the two preceding tools in tandem.  The tailored complement
places a minimizing divisor in a bounded chart; integral approximation
then replaces it by a prime order of small discrepancy.  This turns a
failure of comparison on the base into a violation of the corresponding
hypothesis upstairs.

\begin{proof}[Proof of Lemma~\ref{lem:A-comparison}]
We induct on the dimension of the total space; the birational case was
settled in Subsection~\ref{subsec:A-mori}.

\emph{Extraction and effectivity.}
Extract using \(a_0\) with \(\delta_*\) fixed small and take the Mori
step \(X_1\to W\).
Effectivity of all parameters on the extraction holds by monotonicity, or in the two-sided case by the ratio hypothesis taking \(\delta_*<\min(\delta,1)/3\), after trimming \(L_i\to\infty\) more slowly.
Indeed concavity using the opposite endpoint of a fixed interval controls interior upper ratios too.
For example, if \(0\le t\le T\), then concavity at zero gives
\[
 a_t(T')\le (1+t/T)a_0(T')-(t/T)a_{-T}(T').
\]
The hypothesis on all test subsequences makes the right side
\((1+o(1))a_0(T')\), uniformly on the extracted divisors.
The other sign is identical.  Apply this for each fixed rational
\(T\), then choose a slowly increasing interval by diagonalization.
At subsequent Mori steps the discrepancy families are discriminants of
the original family.  They need not have a common global piecewise
affine presentation.  For each fixed prime \(P\), however, resolve the
original family over its generic DVR and include the fibre in the snc
support.  In the trimmed range, original divisorial valuations trivial
on the current function field already have positive discrepancy.
The discriminant is therefore the minimum of the finitely many vertical
component ratios on that resolution.  Minima compose under the nested
field restrictions, so the finite rational parameter transitions used
below persist at each step.  Generic klt over the next base can likewise
be tested on a resolution of the original family over its generic field.

\emph{Generic klt over the next base.}
We claim generic klt over \(W\) in the Mori step for all bounded parameters eventually.
Otherwise take a sign and a first horizontal lc threshold \(t_*\ne0\) with \(|t_*|\) bounded, and a horizontal minimizing divisorial \(T\).
This first threshold is rational by resolution; if tested on the original family, a place giving it restricts nontrivially to the current field (by the previous generic klt), hence gives an lc place there too.
Tailor the chart using \(a_{t_*}\).
Concavity at this place gives
\[
              a_{\pm L}\le (1-L/|t_*|)a_0
\]
for arbitrarily large fixed rational \(L\), sign as chosen.
Let \(w\) be the chart vector representing \(T\).
Put \(u=w/a_0(w)\), and bound \(\lvert t_*\rvert\) by a fixed
\(T_0\).  If \(\lvert u\rvert\) were bounded along a subsequence,
the uniform homogeneous Lipschitz bound would bound
\(\lvert a_{\pm L}(u)\rvert\) independently of fixed \(L\), whereas
\(a_{\pm L}(u)\le1-L/T_0\).  Choosing \(L\) large gives a
contradiction.  Thus, after a further subsequence,
\(\lvert u\rvert\to\infty\).
Lemma~\ref{lem:A-lattice-approximation}\textup{(i)} gives
\(\lvert n_i-s_i u_i\rvert\le C_*s_i\), \(n_i\ne0\), and hence
\[
 0<a_0(n_i)\le (1+CC_*)s_i,\qquad
 a_{\pm L}(n_i)\le(1-L/T_0+CC_*)s_i .
\]
The same \(C\) works for every effective CY discrepancy
\(a_{\pm L}\).
Choose \(L>T_0(1+CC_*)\) after fixing the subsequence.
The restriction of \(n_i\) is \(k_i\ord_{T_i}\), with
\(k_i\ge1\).  Therefore \(a_0(T_i)\to0\), while
\(a_{\pm L}(T_i)<0\), contradicting the appropriate ratio
hypothesis.
This proves the claim; trim again to keep generic klt throughout and induce effective generalized data on \(W\).

\emph{Transfer of the ratio hypothesis.}
We must pass the corresponding ratio hypothesis to \(b_t(P)\) on \(W\).
Suppose there is failure with \(b_0(P)\to0\), respectively with \(b_0(P)<\delta'\) where \(\delta'>0\) is fixed sufficiently small below.
Concavity (also using the other sign in the two-sided case) gives a bounded parameter with \(b_t\le0\).
Move from zero along a sign to the first \(b_t=b_0/2\); we may take \(b_s\le2b_0\) up to that time: if there is a higher peak use the opposite sign, where the drop is then forced.
In the one-sided case this uses monotonicity instead.

There is a rational parameter \(p\) on this segment and a simultaneous minimizer over \(P\) with
\[
        a_0(T)/e(T)\to0\quad
        (\text{respectively }\ a_0(T)/e(T)\le 10\delta'),\qquad
        a_p(T)\le \tfrac12 a_0(T).
\]
Indeed take a cutoff radius \(\sqrt{b_0}\) or \(8\delta'\).
At zero all minimizers have ratio \(a_0/e=b_0\).
At the half-drop any minimizer below the cutoff works.
If minimizers exit the cutoff first, at a transition interpolate where a minimizer below and one at/above tie, to obtain a ratio near the cutoff, where the desired discrepancy drop holds since \(a_p/e=b_p\le 2b_0\).
It suffices to track the finite minimizing candidates from log resolution (or the divisorial restrictions of the finite candidates minimizing in the original family at each parameter; restriction of any minimizing original candidate gives a current minimizer by composition); use the LC-place connected interpolation above at the rational transition parameter.

Tailor the vertical chart to this \(T\).
In its notation \(\alpha=\rho\,a_0(T)/e(T)\) and concavity gives a decrease as above by arbitrarily large factors at a fixed signed \(L\).
If \(\alpha\to0\), kernel approximation and uniform Lipschitz at fixed uniformizer value give again
\[
                   a_0(n,E)=O(s),\qquad a_{\pm L}(n,E)<0
\]
for sufficiently large fixed \(L\), impossible.
Otherwise we are in the two-sided case with \(\alpha\le C\delta'\) but bounded below after passage to a subsequence.
Use the bounded Dirichlet approximation with \(\gamma\) so small that the Lipschitz error for \(a_0\) is \(<\delta/2\).
The constants here (also ramification and \(J\)) are uniform for this Mori step, so choose \(\delta'\) small enough that \(E\alpha<\delta/2\).
A sufficiently large signed \(L\) fixed \emph{on the subsequence} still makes the discrepancy negative by the uniform Lipschitz bound even for \(a_{\pm L}\).
This contradicts the strong ratio hypothesis.
We have thus transferred that hypothesis with fixed \(\delta'>0\) (and the one-sided hypothesis without a fixed cutoff as required).
Iterate using \(\dim W<\dim X\), stopping if birational over \(S\).
This proves Lemma~\ref{lem:A-comparison}.

\end{proof}

\begin{proof}[Proof of Lemma~\ref{lem:A-pullback-phases}] Absorb the input multiple into \(H\).
Use the same Mori and discriminant steps with fixed klt CY \(a\).
For \(b(P)\to0\) at a step choose a minimizer, so in the vertical chart \(\alpha=a(w,e_*)/e_*=\rho b(P)\to0\).
Here write \(h=v_P(H_W)\) where \(H_W\) is the base pullback.
Use the kernel argument on the augmented torus with sampling vector \(l(U,\rho h)\).
Write \((u_*,z_*)\) for an extension of the signed limits divided by \(l\alpha\) on the bounded-frequency subgroup \(\Gamma\).
The same sampling equidistributes modulo the augmented lattice after subtracting \(s(u_*,z_*)\).

Every vector in the Lie space of \(\Gamma^\perp\) must have last coordinate zero.
Otherwise choose its sign with negative last coordinate and target a small multiple \(\xi v\) with \(\xi\to0^+\) on that torus.
By taking intervals \(0<s\ll\xi\to0\) sufficiently slowly and approximating the target to \(o(\xi)\) along a diagonal subsequence, the integral matched valuation has angular error \(O(\xi)\), hence \(a(n,E)=O(\xi)\), but signed phase \(\xi v_{\rm last}+o(\xi)\), contradicting the lower sign bound.
The sign bound applies to integral multiples whenever their total
discrepancy tends to zero.  Indeed, for \(v_i=k_i\ord_{T_i}\) with
\(a(v_i)\to0\), the primitive discrepancies also tend to zero.
If their phase lifts are
\(r_i=O(a(T_i))\), \(r_i\ge-\varepsilon_i a(T_i)\), with
\(\varepsilon_i\to0\), then \(k_i r_i\) are lifts with
\[
 k_i r_i=O(a(v_i)),\qquad
 k_i r_i\ge-\varepsilon_i a(v_i).
\]
They tend to zero, so no phase wrapping affects the sign.
The projection of the compact subgroup \(\Gamma^\perp\) to its
last circle now has zero Lie algebra, hence finite image.
Some \(q>0\) annihilates this image; by character duality
\((0,q)\in\Gamma\).  Consequently
\({\rm dist}(q l\rho h,\mathbb Z)=O(l\alpha)\).
Its normalized residual limit is nonnegative: indeed \(z_*\ge0\) by the same sign test targeting zero on the torus with precision \(o(s)\), angular error \(O(s)\).
Pass to a subsequence with \(l,\rho\) fixed.

For every test sequence \(b(P)\to0\) we thus have a subsequence and a fixed multiple there giving both the \(O(b)\) bound and the lower \(-o(b)\).
There exists a \emph{single} fixed multiple giving the upper size bound
for all such sequences.  Otherwise choose diagonal violations with
\[
 {\rm dist}(k!h,\Z)>k!\,k\,b(P),\qquad b(P)\longrightarrow0.
\]
The subsequence property gives a fixed good multiple \(M\) on a
further subsequence, with
\({\rm dist}(Mh,\Z)\le Cb(P)\).
For large \(k\), \(M\mid k!\), and thus
\[
 {\rm dist}(k!h,\Z)
 \le (k!/M){\rm dist}(Mh,\Z)
 \le (k!/M)C b(P),
\]
a contradiction.
Fix such a multiple \(M_0\).  If its small lift has a negative
ratio to \(b(P)\) bounded away from zero along a subsequence, apply
the subsequence property again and take a common multiple of \(M_0\)
and the resulting good multiple.  Both small lifts, multiplied by
fixed positive integers, tend to zero and represent the same phase,
so eventually they are equal.  Their incompatible lower signs
give a contradiction.  This proves the required lower
\(-o(b(P))\) bound for \(M_0\).
This gives the required induction and proves Lemma~\ref{lem:A-pullback-phases}.

\end{proof}

\begin{proof}[Proof of Lemma~\ref{lem:A-linear-chamber}] Use the Mori step extracting by \(a+t_0d\).
Since \(a\) has bounded denominator and \(d\ge0\), all extracted divisors with this discrepancy tiny have \(a=0,\ d\) tiny.
Thus in the extra-data case too, effectivity on the extraction holds uniformly, by the upper bounds on \(a=0\).

\begin{claim}[Quantized chart cells]\label{claim:A-quantized-cells}
The bounded-denominator functions \(a,h\), when present, have only a
bounded finite set of possible slopes in a horizontal chart, or in
integral kernel coordinates on a vertical slice \(e_*=1\).
Their vertical affine constant terms belong to a rational grid with
bounded denominator, although their magnitudes need not be bounded.
Their common cells are cut out by a fixed finite set of possible
normal vectors, with offsets of bounded denominator.  They are
pointed, their vertices have bounded denominator, and they satisfy a
uniform polyhedral distance bound.
\end{claim}

\begin{proof}[Proof of the claim]
Horizontally, the functions of bounded denominator \(a,h\) have homogeneous linear pieces with slopes in a fixed bounded finite rational grid: their Lipschitz constants are bounded using values of effective CY discrepancies at finitely many fixed parameters, and their values on integral vectors lie in a prescribed rational grid, which controls the linear form on any full-dimensional open conical piece (take lattice differences there).
Thus their cells can be cut by the comparisons of those finitely many linear forms.
Vertically work on \(e_*=1\), writing \(x=x_0+U\) as before in an integral kernel basis.
Linear parts in \(U\) again belong to a bounded finite rational grid, and the possibly unbounded constant terms at \(x_0\) to a rational grid with bounded denominator (test lattice differences on the open \emph{homogeneous} pieces, using integrality after multiplying \(x_0\) by \(l\)).
Coordinates forced zero on the slice can first be dropped (use \(x_0\) for the reduced system), giving full relative dimension in \(x_0+\ker M\).
More explicitly, an interior of linearity there lifts by homogeneity to an open in the linear space of matches.
Take a lattice point interior to this homogeneous piece and translate a sufficiently large positive multiple of it by each of \((l x_0,l)\) and the kernel basis vectors (with last coordinate zero).
Both evaluations in each difference are on integral matches in the same linearity piece, which gives the intercept and kernel-slope quantization regardless of the size of \(x_0\).
Closures of cells can be cut by inequalities from a finite set of normals, with offsets of bounded denominator: compare pairs of the affine forms occurring and include nonnegativity constraints.
Values are affine on these comparison cells by continuity.
Vertices have bounded denominator, and the polyhedra are pointed.
We also use the elementary polyhedral error bound (Hoffman's bound): distance to a nonempty polyhedron with rows in a fixed finite list is bounded by a constant times the maximum violation.
Here is the distance estimate explicitly.  Write a nonempty polyhedron
as \(P=\{y:R_\nu y\le b_\nu\}\), with rows \(R_\nu\) in the
fixed finite list.  If \(y\) is the nearest point of \(P\) to \(x\),
express \(x-y\) in its normal cone using independent active rows:
\[
 x-y=\sum_{\nu\in I}\lambda_\nu R_\nu^{\mathsf T},
 \qquad \lambda_\nu\ge0,\qquad
 \sum_{\nu\in I}\lambda_\nu\le C\lvert x-y\rvert .
\]
The constant is uniform because only finitely many independent row
sets can occur.  Taking the scalar product with \(x-y\) gives
\[
 \lvert x-y\rvert^2
   =\sum_{\nu\in I}\lambda_\nu(R_\nu x-b_\nu)
   \le C\lvert x-y\rvert
      \max_\nu(R_\nu x-b_\nu)_+.
\]
Cancel \(\lvert x-y\rvert\) when it is nonzero.
\end{proof}

\emph{Horizontal signs and generic klt.}
First in the extra-data case, \(h<0\) on a horizontal \(a\)-lc place over \(W\) is impossible eventually.
Tailor using \(a\).  The homogeneous comparison cells for \(a,h\)
come from finitely many rational cones.  On a cell containing a
witness, the face \(a=0\) is rational and \(h\) is linear.
Some extremal ray of that face has \(h<0\).  The possible primitive
integral ray generators form a bounded finite set, so there is a
bounded nonzero integral witness \(a=0,\ h<0\).
But \(c\) on it is bounded by Lipschitz (again by fixed parameters), and \(h\) quantized and \(\ge-\eta_i c\), impossible.
Next generic klt for \(a+\lambda(h+\theta c)\) holds uniformly for small \(\lambda>0\), \(0<\theta\le\theta_0\): otherwise take \(\lambda_i\to0\) equal to a first positive lc threshold with the corresponding fixed \(\theta_i\) in each example.
There is such a first threshold by common resolution, since on horizontal \(a=0\) we have \(h\ge0,\ c>0\).
Tailor using this threshold place.
In the homogeneous chart distance from its vector \(w\) to \(a=0\) (including the zero vector) is bounded by \(C_1 a(w)\) by the slope/polyhedral bound.
On that set \(h+\theta_i c\ge0\), by the divisorial conditions and rational approximation there.
Thus \(h+\theta_i c\ge-C_2 a(w)\) at \(w\), inconsistent with the threshold tending to zero (\(a(w)>0\)).

\emph{Lexicographic minima and their denominators.}
For \(P\) over \(W\) denote the claimed lexicographic minima by \(a',d'\) (and \(a',h',c'\)).
They exist: on a common resolution over the DVR use the vertical components and the formula for discrepancies by monomial orders plus the nonnegative reduced-log discrepancy.
Horizontal \(a\) is nonnegative, and in the extra-data case on horizontal \(a=0\), \(h\) is nonnegative and \(c>0\).
Explicitly include the fibre among the divisors on the resolution, so \(e\) is a nonnegative linear combination of vertical orders only.
At a generically klt parameter a minimizer has reduced-log discrepancy zero (on the resolution) and all its positive orders in those coordinates lie on minimizing vertical components, by subtracting the minimum times \(e\); thus tracking these finitely many components captures departures from the faces in the argument.
For each example and \(P\), sufficiently small strictly positive parameters in lexicographic order thus have only lexicographically minimizing components, and give a simultaneous minimizer for tailoring using a generically klt parameter.
On that chart the minimum \(a/e_*\) is \(\rho a'\), and likewise for \(h\) on the minimizing locus if used.
By linear optimization at cell vertices on the slice \(e_*=1\), both of these minima have bounded denominator.
Thus so do \(a'\) and \(h'\), uniformly in \(P\).

\emph{Uniformity of the chambers.}
If minimizers of \(a+t d\) leave the \(a\)-minimum for some \(t\) tending to zero, take the first (strictly positive rational) transition where minimizers of both types tie.
Connected interpolation gives a simultaneous minimizing place with \((a/e)-a'(P)\) arbitrarily small positive.
Tailor at that place for the transition discrepancy.
On the cell of its normalized weights (value \(e_*=1\)), the affine
function \(a-\rho a'(P)\) is nonnegative.  Its minimum is attained at
a vertex, because the cell is pointed.  Both the vertex values of
\(a\) and \(\rho a'(P)\) lie in a fixed rational grid.  A strictly
positive minimum would therefore be bounded below by a fixed positive
number.  Taking the interpolated excess smaller than this number forces
the face \(a=\rho a'(P)\) in that cell to be nonempty.
Let \(\varepsilon=a(w)-\rho a'(P)>0\).
The distance estimate in Claim~\ref{claim:A-quantized-cells}
gives \(w_0\) on that face with
\(\lvert w-w_0\rvert\le C_1\varepsilon\).
If \(C_2\) is the Lipschitz constant of \(d\), then
\[
 (a+td)(w_0)-(a+td)(w)
 \le-\varepsilon+tC_1C_2\varepsilon<0
\]
for a uniformly small \(t>0\).  This contradicts minimality;
continuity permits the comparison at real weights and rational
approximation recovers divisorial tests.
This proves uniform exact minimization on the \(a\)-face for \(a+t d\).
Identical reasoning with \(h+\theta c\) yields uniform small \(\lambda\) with \emph{all} minimizers still on the \(a\)-face for every \(0<\theta\le\theta_0\) in the extra-data case.
Fix a uniformly positive small such \(\lambda\).
If first positive \(\theta\)-transitions away from the \(h\)-minimum within the \(a\)-minimum approach zero, now interpolate with tiny positive excess of \(h/e\) above \(h'\); all simultaneous minimizers stay on the \(a\)-minimum.
Use joint \(a,h\) cells with the additional affine face \(a=\rho a'\); the grid and distance argument within this face applied to \(h\) contradicts minimality for \(h+\theta c\).
Within the \(a\)-minimum the minimizer condition no longer depends on the magnitude of \(\lambda>0\), only on \(\theta\).
Thus the full formulas hold uniformly on a smaller sector.

\emph{Induction on the base.}
For induction \(a'+t d'\) gives klt data for positive \(t\) by the minimum formula, and data at zero by the endpoint limit already explained.
In the extra-data case on \(a'=0\), \(d'\le c'\), \(h'\ge-\eta_i c'\), and testing at a \(d'\)-minimizer at fixed uniform small \(\theta>0\) gives \(h'+\theta c'\le C_3 d'\) by the exact formula.
Hence \(c'\le C_4 d'\) and \(h'\le C_3 d'\) there.
All hypotheses now hold at the next step.
Any negative \(h\) at \(a=0\) vertical in this step marks a negative \(h'\) at \(a'=0\), by restriction.
Induct to the base to obtain all assertions. \end{proof}
